%% P05 \dash{} Inner product, norms and projection
%%
%% Stub. The brief below is the contract for this program: what it must argue,
%% and what it must buy the reader. Writing the program means deleting the
%% \programstub{} block and replacing it with the program. If the brief turns
%% out to be wrong once the mathematics has been worked, change it and record
%% the contradiction in CLAUDE.md under *Resolved questions*.

\program{Inner product, norms and projection}
\label{prog:P05}

\programstub{%
Argument: an inner product is the only thing that gives a vector space
angles and lengths; a projection is the closest point in a subspace;
different norms measure different things and disagree about which vector is
bigger. Payoff: dot product, cosine and Euclidean distance compared on the
same data, with the case where they rank differently worked out \dash{} the
thing a vector database forces you to choose and nobody explains; why
normalising embeddings makes cosine and dot product the same query and what
you lose; L1 against L2 as a shape rather than a slogan, and why the L1
ball's corners are the whole of the sparsity argument. Also the fact the
book cashes in twice later: in high dimension, two random vectors are almost
always nearly orthogonal, measured rather than asserted.

\medskip
\textbf{Planned length:} 60 frames.
}
